\section{Challenges Faced and Future Recommendations}

\subsection{Challenges Faced}
Several technical and analytical challenges were encountered in this project:
\begin{itemize}
  \item Post-level engagement metrics do not directly measure real investment execution or financial outcomes.
  \item Text content is unstructured and noisy, requiring substantial preprocessing before advanced NLP.
  \item Engagement variables are heavily skewed, making naive summary statistics sensitive to extreme viral posts.
  \item Feature-stage performance can appear near-saturated because the virality target is score-derived, which raises leakage-like ease concerns.
  \item The 2020--2021 window reflects an unusual market regime, limiting generalization to other periods.
\end{itemize}

\subsection{Future Recommendations}
Future extensions that would improve explanatory power include:
\begin{itemize}
  \item Replace score-derived virality labels with forward-looking targets (e.g., future engagement windows) to reduce label leakage risk.
  \item Strengthen text modeling using contextual embeddings and sarcasm-aware sentiment methods tailored to finance slang.
  \item Integrate external market indicators (price, volume, volatility) for cross-domain linkage between discourse and market movement.
  \item Use time-aware validation (rolling or forward-chaining splits) to better reflect deployment realism.
  \item Extend cross-community comparisons beyond WSB to test transferability of findings.
\end{itemize}

These extensions would strengthen causal interpretation, improve deployment realism, and increase generalizability of investor-attention analysis.

Evidence context:
\begin{itemize}
  \item Model diagnostics and limitations: \verb|.../artifacts/json/05_feature_ablation_table.json|
  \item EDA anomaly and dependence diagnostics: \verb|.../artifacts/json/03_eda_advanced_stats.json|
\end{itemize}
